\documentclass[11pt]{article}
\usepackage[utf8]{inputenc}
\usepackage{amsmath, amssymb}
\usepackage{booktabs}
\usepackage{geometry}
\geometry{margin=1in}
\usepackage{hyperref}

\title{From Free--Energy Loops to Human Values:\\A Hub--Centric Precision Model}
\author{Gunnar~Zarncke\\AE~Studio}
\date{July~2025}

\begin{document}
\maketitle

\begin{abstract}
The preceding paper catalogued thirteen free--energy--minimising loops in the vertebrate brain.  Each loop closes a prediction--error feedback and projects into one or more neuromodulatory or salience ``hubs'' that broadcast precision weights.  Here we introduce the \emph{loop--hub--value} (LHV) model: an architectural proposal in which a small set of hubs transforms high--dimensional error vectors into low--bandwidth scalars that are then read out as human value tags.
We (i) formalise the two transfer steps (\emph{loop~$\rightarrow$~hub} and \emph{hub~$\rightarrow$~value}); (ii) compile an empirical mapping between hubs and intrinsic values, drawing on the annotated ``brain.py'' resource; and (iii) outline falsifiable predictions.
Details of agent implementation and learning dynamics will appear in Part~III.
\end{abstract}

\section{Motivation}
Biological agents cannot transmit every prediction error upstream; long--range bandwidth is constrained by wiring and metabolic costs.  Evolution appears to have solved this with \emph{hub bottlenecks}: mid--brain nuclei and limbic structures that \emph{compress} multidimensional errors into one or two slow neuromodulatory signals.  Cortex, basal ganglia and cerebellum decode those signals as \emph{values} that steer policy selection.

\section{Formal Skeleton of the LHV Model}
\paragraph{Error compression.}  Each loop $i$ yields an error vector $\boldsymbol{\epsilon}_i(t) \in \mathbb{R}^{d_i}$.  A hub $h$ receives a weighted sum over incoming loops
\begin{equation}
    s_h(t)\;=\;\sigma\Bigl(\sum_{i \in \mathcal{I}(h)} w_{ih}\,\lVert\boldsymbol{\epsilon}_i(t)\rVert_{p}\Bigr),
\end{equation}
where $\sigma$ is a saturating non--linearity and $w_{ih}$ are learned transfer weights.
\paragraph{Value tagging.}  A cortical decoder $D$ maps hub signals to symbolic values
$\mathcal{V} = \{v_1,\dots,v_m\}$ via a softmax
\begin{equation}
    P\bigl(v_k\mid \{s_h\}_h\bigr)\;=\; \mathrm{softmax}_k\bigl(\beta_k + \sum_h \alpha_{kh}\,s_h\bigr),
\end{equation}
with $\alpha_{kh}$ contextually gated by attention and task set.

\section{Empirical Hub--Value Map}
Table~\ref{tab:hv} aggregates
(1) the thirteen loops from Revision~2; (2) the hubs that broadcast their precision; and (3) the intrinsic values manually assigned to hubs in the open--source script \texttt{brain.py}.  Region--value attributions follow the dictionary in that script. ([raw.githubusercontent.com](https://raw.githubusercontent.com/GunnarZarncke/brain-to-values/main/brain.py))

\begin{table}[h]
\centering
\footnotesize
\begin{tabular}{@{}lll@{}}
\toprule
Hub & Primary loops feeding hub & Example intrinsic values \\
\midrule
Periaqueductal~Gray (PAG) & Threat\/pain; Pavlovian threat tagging & Protection, Non--suffering, Freedom\\
Ventral~Tegmental~Area (VTA) & Curiosity; Valuation regret & Learning, Diversity, Legacy\\
Locus~Coeruleus (LC) & Precision\/timing & Achievement\\
Nucleus~Accumbens (NAcc) & Valuation\/agency & Happiness\\
Orbitofrontal~Cortex (OFC) & Policy--regret valuation & Happiness, Appraisal\\
Hypothalamus & Metabolic; Steroidal timing & Longevity, Caring\\
HPG~axis & Steroidal timing & Reputation\\
Insula & Metabolic; Morphological repair & Fairness, Purity\\
Anterior Cingulate Cortex (ACC) & Coalition entropy & Justice, Respect\\
Amygdala & Pavlovian threat tagging & Virtue\\
Septal~Nuclei & Reproductive physiology & Loyalty\\
Superior~Colliculus & Perimeter defence (external) & Nature, Truth\\
Primary Visual Cortex & Scene novelty loop$^{\dagger}$ & Beauty\\
\bottomrule
\end{tabular}
\caption{Hub--value correspondences.  $^{\dagger}$The scene novelty loop is hypothesised and will be detailed in Part~III. Values sourced from \texttt{brain.py}.}
\label{tab:hv}
\end{table}

\section{Predictions}
\begin{enumerate}
  \item Disrupting a hub (chemogenetically or by deep brain stimulation) should \emph{simultaneously} perturb all values anchored to that hub, even when the upstream loops are intact.
  \item Task contexts that raise a loop's error norm $\lVert \boldsymbol{\epsilon}_i \rVert$ will modulate subjective value ratings in proportion to the learned weight $w_{ih}$.
  \item Artificial agents endowed with an LHV bottleneck will exhibit more interpretable preference shifts when hub weights are manipulated than agents with direct reward summation.
\end{enumerate}

\section{Outlook}
Part~III will instantiate LHV in a hierarchical active--inference agent and benchmark alignment transparency.  Additional work is required to (i) learn $w_{ih}$ and $\alpha_{kh}$ from data, (ii) integrate steradian bandwidth constraints, and (iii) test cross--species scalability.

\section*{Acknowledgements}
The value annotations stem from the open GitHub resource at \url{https://github.com/GunnarZarncke/brain-to-values}.  I thank the contributors for maintaining clear documentation.

\end{document}
